% !TeX root = ../papers/nc-p3.tex
\section{中央二重の有限点窓の唯一の曲線}
\label{sec:double-middle-finite-windows}
\begin{theorem}\label{thm:double-middle-finite-windows}
有限点グラフ上の $0,1,1,2$ 型の全正則窓は
$\alpha\in k\setminus\{0,1,-1\}$ による一族で、添字保持同型は
$\alpha'=\alpha$ または $\alpha'=\alpha^{-1}$ と同値である。全て最小周期1で多重安定、
点スキームは長さ9の二点と二つの単純点からなる。$\alpha=-1$ の同じ表示も正則で点対応は正次元である。
従って定理\ref{thm:simple-weight-finite-windows}と本定理で、有限点グラフと非自明トーラスを同時に持つ
全窓のパラメータ・例外値・印なし同型・最小周期が決まる。いずれの射影安定群の恒等成分も乗法群である。
\end{theorem}
\begin{proof}[中央四文字と二つの鉛筆]
$x,U_i,z$ を重み別に取る。全体重み4のうち両極だけの成分は交互に並べた二語のみで、どちらも非零である。
四位置の $x,z$ の倍率を使って両係数を1にする。中央四文字成分を $P$ とし、$xz,zx$ を位置 $i,i+1$ に持つ成分の
残り二文字の行列をそれぞれ $A_i,B_i$ とする。$x,z$ を位置 $i,i+2$ に持つ成分の行列を
行 $i+1$、列 $i+3$ の順に $C_i$ とする。これらは全て可逆である。実際、重み1対3の切断の両側で
二つの射影が同型であり、各行列はその切断の一隅である。

重み2対2のSchur補は次の二式である。反対位置の二関係から来る中央行列は
$A_{i+2},B_{i+2}$ そのものである。
\[
 P=A_{i+2}\otimes A_i+B_{i+2}\otimes B_i\qquad (i=0,1).
\]
各辺の二行列は二次元鉛筆を張る。そうでなければ中央二直線の積に曲線が残る。
また鉛筆は共通の行核・列核を持たない。各鉛筆と階数1の二次曲面の交叉は、
相異なる二点または重複した一点である。

第一鉛筆が二点型なら左二位置の基底で $00,11$ を張る形となる。
$P=00\otimes H+11\otimes K$ の次切断の階数は $\rk H+\rk K=2$。
両者の行直線と列直線はそれぞれ独立なので $P=0000+1111$ にできる。
重根型なら第一鉛筆は $00,01+10$ を張り、$P=00\otimes H+(01+10)\otimes K$ となる。
次切断は転置や行列並べ替えを除き
$\left(\begin{smallmatrix}H&K\\K&0\end{smallmatrix}\right)$ の階数2を要求する。
従って $K$ は階数1で $H$ はそのSegre接平面にあり、最後の二位置で $K=00,H=01+10$ にできる
（行列鉛筆の正則性で両交差係数が非零である）。よって $P=1000+0100+0010+0001$。
これで全辺の型がそれぞれ同一になった。
\end{proof}

\begin{proof}[混合切断の一巡]
二点型では鉛筆の順序付き基底を $00,11$、重根型では $00,01+10$ とする。
$D_i$ をこの基底での $A_i,B_i$ の二列とすると
$D_{i+2}=J D_i^{-\mathsf t}$ であり、$J$ は二点型で恒等、重根型で二座標の交換である。
さらに、重み1対3の行を $xU,Ux$、列を $zU,Uz$ とすると
\[
 \begin{pmatrix} C_i&B_{i+3}\\ A_{i+1}^{\mathsf t}&C_{i+1}^{\mathsf t}\end{pmatrix}
 \quad\Longrightarrow\quad C_{i+1}=B_{i+3}^{\mathsf t} C_i^{-\mathsf t}A_{i+1}.
 \tag{1}
\]
従って $D_0,D_1,C=C_0$ が四回の式(1)で閉じることだけが残る。

二点型では中央文字の対角倍率と極文字の倍率を使い
\[
 D_0=\begin{pmatrix}1&1\\1&\alpha\end{pmatrix},\qquad
 D_1=\begin{pmatrix}1&1\\1&\beta\end{pmatrix}
\]
にできる。$\alpha,\beta\ne0,1$ である。実際中央二文字の倍率比 $r_i$ は
$\prod r_i=1$ としてよく、$r_2r_3,r_3r_0$ を独立に指定できる。
各列には極文字の倍率積 $x_0z_1,z_0x_1,x_1z_2,z_1x_2$ が独立に作用する
（残りの二つの極倍率は交互極語の正規化から決まる）。
(1)を四回用いると $C=\operatorname{diag}(l,m)C\operatorname{diag}(m,l)$ となる。ただし
\[
 l=(\beta-1)/(\alpha-1),\qquad m=\alpha l/\beta .
\]
可逆 $C$ があるには $lm=1$ または $l^2=m^2=1$ が必要である。
前者は $(\beta-\alpha)(\alpha\beta-1)=0$、後者でも $l=1$ または $m=1$ なら $\beta=\alpha$、
両方 $-1$ は不可能である。従って $\beta=\alpha$ で任意の $C$、または $\beta=1/\alpha$ で
対角 $C$ だけが残る（重複時は前者）。後者では共通倍率および両側からの同じ対角合同で $C=I$ にできる。

重根型の中央基底の残余は、位置別のスカラーを除き
$e_0\mapsto e_0,e_1\mapsto k e_1+m_i e_0$, $\sum m_i=0$ である。
$P$ の線を保つには $k$ が共通であり、これらは全てその線を保つ。
各列の第二座標は可逆性から非零で、極倍率を使い第一・第二鉛筆で1に揃える。
第一座標は $m_2+m_3,m_3+m_0$ の平行移動と共通の $k^{-1}$ 倍を使って
\[
 D_0=\begin{pmatrix}0&-1\\1&1\end{pmatrix},\qquad
 D_1=\begin{pmatrix}0&-\beta\\1&1\end{pmatrix},\qquad \beta\ne0
\]
にできる。$E=E_{01}$ とすれば一巡条件は
$C=\beta^2(I+(\beta-1)E) C(I-(\beta-1)E^{\mathsf t})$。
右の単冪部分の固有値は1なので、$\beta=1$ で任意の $C$ か、
$\beta=-1$ で $EC=CE^{\mathsf t}$ となる。後者は
$C=\left(\begin{smallmatrix}u&v\\v&0\end{smallmatrix}\right)$, $v\ne0$ で、
残余平行移動 $m_0=m_2=-m_1=-m_3$ と共通倍率で交換行列にできる。
\end{proof}

\begin{proof}[二つの二元三次式と同型]
辺の最低文字を双方1にする。一方だけ非零なら混合グラフと $xz$ の式から他方の全座標が零になり、
最高文字でも同じである。中央二座標を $u,v$ とすると重み1より
$v=-C_i^{-\mathsf t}A_{i+1}u$。重み2で両最高座標は $-u^{\mathsf t}B_{i+2}v,-u^{\mathsf t}A_{i+2}v$ となる。
重み3の二関係は従って次の二元三次式の同時消滅である。
\[
 (u^{\mathsf t}B_{i+2}v)C_{i+2}v+(u^{\mathsf t}A_{i+2}v)B_{i+1}u=0.
\]
非零共通根があれば非自明な軌道がある。
二点型 $\beta=\alpha$ と重根型 $\beta=1$ は一般行列 $C$ の関数体上で始点0の二式に共通二次因子がある。
重根型 $\beta=-1$ は共通一次因子がある。終結式が恒等零なので可逆な全特殊値でも有限性は成立しない。
残る二点相反型では、四位置と極反転後の計八終結式は全て $\alpha+1$ の二乗に
$\alpha,\alpha-1$ の単元を掛けたものになる。中央だけの辺は二点型鉛筆の二つの単純点のみである。
正確な二次行列反復・共通因子・Sylvester行列式は \path{computations/verify_double_middle_torus.py} で検算できる。
これで定理のパラメータ領域と復元式が得られる。逆向きの正則性は命題\ref{prop:double-middle-finite-pbw}による。
$\alpha=-1$ でもその証明の可逆性と三文字の基底は保たれる。

二点型の二つのSegre分解項は順序を除き一意であり、各位置で交換するなら四位置全てで同時に交換する。
列倍率と行倍率を除く $D_0$ の交差比が $\alpha$ で、同時交換は $\alpha^{-1}$ を与える。
極文字の交換は $A_i,B_i$ を交換するので同じ反転しか加えない。相反型の対角 $C$ は常に一類であるから十分性も従う。
有限グラフでは極大トーラスの共役によりこれが印なし同型であり、一巡回は $\alpha$ を $\beta=\alpha^{-1}$ に移すため周期1となる。

最後の標準形の支持は二つの交互極語、中央のみで全て同じ文字の二語、そして任意の異なる二位置に $x,z$ を置き
残る二位置で同じ中央文字を使う24語である。それぞれの群での一様平均を合わせると全位置四文字の頻度を均一にでき、
一位置だけの相違はない。従って前節と同じモーメント判定で多重安定となる。
両節の閉軌道の射影安定群は還元的で階数1である。もし半単純因子があれば、閉点の有限集合を
その連結群が各点ごとに固定する。他方その根トーラスの最低点は単純な最低文字で、
非自明な $\mathrm{SL}_2$ 表現の極重み直線であるから群全体には固定されない。これで恒等成分の主張も従う。
\end{proof}
